\question[3]
In einem lokalen Netzwerk unterscheidet man die Komponenten \emph{Endgerät}, 
\emph{Verbindung} und \emph{Verteiler}.
Erkläre für jede dieser drei Komponenten kurz ihre Aufgabe im Netzwerk  (z.\,B. aus deinem eigenen Zuhause oder der Schule).

\begin{solutionbox}{5cm}
 
Endgerät: Sendet und empfängt Daten, ist der eigentliche
Kommunikationspartner (z. B. PC, Laptop, Smartphone, Drucker).

Verbindung: Überträgt die Daten physisch (oder per Funk)
zwischen den Geräten (Netzwerkkabel, Glasfaserkabel, WLAN-Verbindung).

Verteiler: Leitet Datenpakete an die richtigen Endgeräte
weiter, verbindet mehrere Geräte im Netz (z. B. Switch, Router,
WLAN-Access-Point, Modem).
 
\end{solutionbox}